% ============================================================
%  第一章
% ============================================================
\section{1977：在哈佛与牛津之间}
\markboth{第一章\ 1977：在哈佛与牛津之间}{第一章\ 1977：在哈佛与牛津之间}

\noindent{\small\color{shuyiGray}对应底本：原 p.2--p.6。}\par\medskip

\srcp{2}1977 年春天，我在哈佛做 Junior Fellow，Michael Atiyah 来哈佛讲学。他讲的是瞬子。我听完很受触动，去和他交谈；他邀我在当年年底访问牛津。

那年年底我到了牛津。整个访问期间，Michael 讲的是四维规范理论里的瞬子。他和 Vladimir Drinfeld、Nigel Hitchin、Yuri Manin 刚刚借助扭量理论与代数几何发现了瞬子的 ADHM 构造。瞬子当时是物理里最令人兴奋的题目之一，很大程度上是因为 Gerard 't Hooft 的工作：他用瞬子解决了核力的“$U(1)$ 问题”。到 1977 年底，我们又从 Albert Schwarz 的一个观察知道，'t Hooft 关于 $U(1)$ 问题的工作可以理解为 Atiyah--Singer 指标定理的一次应用。

\srcp{3}那个月的大部分时间，我们都在热烈地讨论瞬子。但到了月底，Atiyah 换了题目。他让我注意两篇我没见过的物理论文，一篇是 Peter Goddard、Jean Nuyts 与 David Olive 的，一篇是 Olive 与 Claus Montonen 的。他说他觉得这里面有很深的东西，我应该去伦敦和 Olive 谈谈。

\srcp{4}Goddard、Nuyts 与 Olive 那篇论文的摘要如下：

\begin{abstractbox}{Goddard, Nuyts \& Olive, \emph{Gauge Theories and Magnetic Charge}（Nuclear Physics B125, 1977, pp.~1--28）}
\en{If the magnetic field for an exact gauge group $H$ (assumed compact and connected) exhibits an inverse square law behaviour at large distances then the generalized magnetic charge, appearing as the coefficient, completely determines the topological quantum number of the solution. When this magnetic charge operator is expressed as a linear combination of mutually commuting generators of $H$, the components are uniquely determined, up to the action of the Weyl group, and have to be weights of a new group $H^0$ which is explicitly constructed out of $H$. The relation between the ``electric'' group $H$ and the ``magnetic'' group $H^0$ is symmetrical in the sense that $(H^0)^0 = H$. The results suggest that $H$ monopoles are $H^0$ multiplets and vice versa and that the true symmetry group is $H \otimes H^0$. In this duality topological and Noether quantum numbers exchange r\^oles rather as in Sine-Gordon theory. A physical possibility is that $H$ and $H^0$ be the colour and weak electromagnetic gauge groups.}

\smallskip
\noindent{\color{srcFrame}\rule{\linewidth}{0.3pt}}\par\smallskip

若一个精确规范群 $H$（假定紧致且连通）的磁场在大距离上呈平方反比律，那么作为其系数的广义磁荷就完全决定了该解的拓扑量子数。当这个磁荷算子写成 $H$ 的彼此交换的生成元的线性组合时，各分量在 Weyl 群作用下唯一确定，并且必须是某个新群 $H^0$ 的权——这个 $H^0$ 可以由 $H$ 显式地构造出来。“电”群 $H$ 与“磁”群 $H^0$ 之间的关系是对称的，意思是 $(H^0)^0 = H$。这些结果表明：$H$ 的单极子是 $H^0$ 的多重态，反之亦然；而真正的对称群是 $H \otimes H^0$。在这一对偶中，拓扑量子数与 Noether 量子数交换角色，颇像 Sine--Gordon 理论里的情形。一个物理上的可能性是：$H$ 与 $H^0$ 分别是色规范群与弱电磁规范群。
\end{abstractbox}

\srcp{5}Montonen 与 Olive 的摘要如下：

\begin{abstractbox}{Montonen \& Olive, \emph{Magnetic Monopoles as Gauge Particles?}（Physics Letters B72, 1977, pp.~117--120）}
\en{We present evidence for the following conjecture: when quantized, the magnetic monopole solutions constructed by 't Hooft and Polyakov, as modified by Prasad, Sommerfield and Bogomolny, form a gauge multiplet with the photon, corresponding to a Lagrangian similar to the original Georgi--Glashow one, but with magnetic replacing electric charge.}

\smallskip
\noindent{\color{srcFrame}\rule{\linewidth}{0.3pt}}\par\smallskip

我们为下述猜想提供证据：量子化之后，'t Hooft 与 Polyakov 构造的磁单极子解（经 Prasad、Sommerfield 与 Bogomolny 修改）与光子组成同一个规范多重态，对应于一个类似于原 Georgi--Glashow 模型的拉氏量，只是以磁荷替换了电电荷。
\end{abstractbox}

\srcp{6}Atiyah 向我解释：Goddard、Nuyts 与 Olive（GNO）所谓的“磁规范群”，其实和数论 Langlands 纲领里的 Langlands 对偶群是同一个东西。我相当确定这个联系是 Atiyah 做出的，但据我所知他从未把它写下来。和当时别的物理学家一样，我从未听说过 Langlands 纲领，也没听过 Robert Langlands 这个名字。但 Atiyah 向我保证：数论的 Langlands 纲领是非常深刻的东西，他认为物理学家必定也在通向某种重要的东西。

\begin{yizhu}[Junior Fellow、twistor、instanton]
\emph{Junior Fellow} 指哈佛学会（Harvard Society of Fellows）的初级研究员，是一个不设教学任务、专供年轻人自由研究的职位，Witten 于 1976--1980 年间任此职。这里保留英文，因为“初级研究员”不能传达其制度含义。\emph{twistor theory} 译“扭量理论”，是 Penrose 在 1960 年代提出的时空几何方案；\emph{instanton} 译“瞬子”，是四维欧氏空间上 Yang--Mills 方程的自对偶解，因在时空局部的“瞬时”存在而得名——这个词的字面意思常被误解为“瞬间的粒子”。
\end{yizhu}

\begin{pizhu}{ADHM：瞬子的线性代数}
Atiyah 与 Drinfeld、Hitchin、Manin 的这篇文章题为 \emph{Construction of Instantons}（Physics Letters A65, 1978, pp.~185--187），通常简称 ADHM。它给出的结论是：$\RR^4$ 上 $SU(N)$ 瞬子（自对偶联络，拓扑荷 $k$）的模空间，等同于一组线性数据
\[
  (B_1,B_2,I,J)\in \End(\CC^k)^{\oplus 2}\oplus \Hom(\CC^N,\CC^k)\oplus \Hom(\CC^k,\CC^N)
\]
在实矩映射方程
\[
  [B_1,B_1^\dagger]+[B_2,B_2^\dagger]+II^\dagger-J^\dagger J=0,\qquad
  [B_1,B_2]+IJ=0
\]
的零点集上，再模去 $U(k)$ 的作用。换句话说，一个非线性椭圆方程的模空间，被写成了一个有限维的线性代数问题（一个超 Kähler 商）。这正是 Witten 后来说的“Atiyah 的方法”：把分析问题翻译成代数几何，再翻译成线性代数。
\end{pizhu}

\begin{pizhu}{$U(1)$ 问题与指标定理}
\emph{$U(1)$ 问题}是 1970 年代量子色动力学的一个困惑：按照手征对称性，九种赝标介子中应当有一个很轻的粒子，但实验上 $\eta'$ 的质量却大得多。't Hooft 在 1976 年指出，瞬子会通过费米子的零模诱导出一个破坏轴 $U(1)$ 的有效相互作用，从而解释 $\eta'$ 之所以不轻；而零模的个数恰由 Dirac 算子的指标给出，也就是由瞬子的拓扑荷给出。Schwarz 的观察是：这一整套论证其实是 Atiyah--Singer 指标定理的一个特例。对 Witten 而言，这件事第一次具体展示了“物理中的拓扑量子数 = 几何中的指标”，而这正是他后来反复使用的手法。
\end{pizhu}

\begin{pizhu}{“他从未把它写下来”}
Atiyah 生前没有为 GNO 对偶群 $=$ Langlands 对偶群这一观察写过专门文章，这条联系长期只以口头传统存在；它在物理文献中被明确引用、并成为一整个研究方向的出发点，要等到 2006 年 Kapustin 与 Witten 关于规范理论与几何 Langlands 的论文。本章末尾 Witten 的那句话之所以有分量，是因为它是这条口头传统的当事人证词：一位数学家在一次非正式的交谈里，把两个学科各自独立发现的同一个结构认了出来。
\end{pizhu}
